% =========================================================
% Final Mid-Semester Review Paper
% =========================================================
\documentclass[acmsmall,screen]{acmart}

\AtBeginDocument{%
  \providecommand\BibTeX{{Bib\TeX}}%
}

\usepackage{booktabs}
\usepackage{array}
\usepackage{float}
\usepackage{tikz}
\usetikzlibrary{arrows.meta,positioning,shapes.geometric}

\title{Knowledge Cutoff Effect on LLM-Generated Literature Reviews in Rapidly Evolving Fields}

\author{Sujal Rajesh Ghagare}
\affiliation{%
  \institution{Indian Institute of Information Technology Allahabad}
  \department{Department of Information Technology}
  \city{Prayagraj}
  \country{India}
}
\email{mse2026015@iiita.ac.in}

\keywords{Large Language Models, Knowledge Cutoff, Literature Review,
Generative AI, Retrieval-Augmented Generation, Citation Reliability,
Hallucination, Scholarly Retrieval, Human-in-the-Loop Verification}

\begin{abstract}
Large Language Models (LLMs) are increasingly used to support
literature-review activities because they can summarize research,
organize information, identify common themes, and generate academic
text rapidly. This, however, can be a serious limitation - the time frame for any information that the model has access to. Research
There is no way for a model's current knowledge to see the publication data if it falls after the time set during training.
Decrease the accuracy and timeliness of an LLM-generated literature review.
This is especially true for those emerging areas such as the impact of this work on artificial intelligence, and security.
Innovation in medicine and biotechnology – where new methods, data, models, and empirical discoveries are emerging.
A recent literature review for have access to a current counter example may miss crucial studies, portray out-dated
Drawn from current approaches, point out already addressed research gaps, or paint too old results.
Furthermore, literature reviews written by LLMs can include inaccurate citations and made-up references.

This review integrates the effects of temporal knowledge restrictions, hallucination, retrieval-
augmented generation, scholarly retrieval, and human verification. It organizes mitigation strategies
into retrieval-based, verification-based, human-based, and hybrid
approaches, compares their strengths and limitations, and identifies
open research gaps. The review states retrieval can provide further accessibility to the recent Evidence, but retrieval alone does not assure source relevance, factuality, or proper citation-to-claim.
Reviews. Dependable literature reviews aided by LLMs, thus, demand contemporary academic references, confirmation,
structured synthesis, and human oversight.


\end{abstract}

\begin{document}
\maketitle

\section{Introduction}

Literature reviews are an essential component of scientific research.
They help researchers understand previous work, compare competing
approaches, identify methodological trends, and determine unresolved
research questions. Traditionally, preparing a literature review
requires researchers to search scholarly databases, screen relevant
publications, read individual studies, extract findings, and synthesize
evidence across multiple sources. The increasing availability of Large
Language Models (LLMs) has introduced new possibilities for assisting
with several of these activities.

LLMs can summarize documents, generate explanations, classify
information, and produce structured academic text. These capabilities
make them attractive as research assistants. However, fluent text
generation does not necessarily imply that the underlying information
is complete, current, or supported by reliable sources. A literature
review is especially sensitive to this distinction because its purpose
is to represent a body of evidence rather than simply produce fluent
text.

A central issue is the temporal boundary of model knowledge. Research
published after the information available to a model may be absent from
its internal knowledge. This can be particularly important in rapidly
One of the fastest maturing areas in machine learning: new papers, datasets, benchmarks, models,
And the experimental results are shown within a few seconds. Fig.~\ref{fig:cutoff}
 shows the fundamental difference model knowledge boundary.
And later publications.

\begin{figure}[t]

\centering

\begin{tikzpicture}[>=Latex,

box/.style={

draw,
 rounded corners,
 align=center,
 minimum height=8mm,

text width=25mm

},

smallbox/.style={

draw,
 rounded corners,
 align=center,
 minimum height=8mm,

text width=31mm

}]

\draw[->] (0,0) -- (9,0);

\node[box] at (1.4,0.7) {Older literature\\available};

\node[box] at (4.3,0.7) {Model knowledge\\boundary};

\node[smallbox] at (7.2,0.7) {Newer literature\\may be missing};

\draw[dashed] (4.3,-0.45) -- (4.3,1.35);

\node[align=center] at (2.0,-0.55) {Training / internal\\knowledge};

\node[align=center] at (6.8,-0.55) {Post-cutoff\\research};

\draw[->,thick] (7.2,-1.25) -- (4.8,-1.25);

\node [smallbox] at (2.4,-1.75)

{External scholarly retrieval\\can provide current evidence};

\end{tikzpicture}



\label{fig:cutoff}

\end{figure}

The second is the quality of citations. Walters and Wilder observed that.
ChatGPT-generated bibliographic citations can contain fabricated
references and substantive errors \cite{walters2023}. Hallucination is
Is also a broader issue in generation of natural language, Ji et al. survey
How to detect and prevent AI-generated content that is not
adequately supported by evidence \cite{ji2023hallucination}. Scientific
Text can be very convincing but still include fabricated data,
As shown by the work of Gao et al. on the scientific produced by ChatGPT
abstracts \cite{gao2023abstracts}.

These problems are of greater importance to a literature review than any one.
Summary. A review necessitates coverage of a topic, comparison across
Studies, interpretation of change over time, and identification
Of research gaps. The absence of up-to-date literature reviews can significantly impact to the productivity of.
Both individual claims and the overall synthesis.

This review was designed to answer the following questions:

\begin{enumerate}

How could information cutoff have an impact on the completeness and

interpretation of LLM-generated literature reviews?

•What retrieval and verification methods can mitigate them

effects?

\item What kind of limitations could be expected if the information were to be given now

to an LLM?

\item Research directions to enable more trusted literature

synthesis in rapidly evolving fields?

\end{enumerate}

The major improvements are a logical timetable analysis,
A taxonomy of mitigation approaches a comparative synthesis of
Retrieval and verification strategies and a practical workflow for
LLM-assisted literature reviews.

\section{Background and Conceptual Foundation}

\subsection{Large Language Models}

A neural language model trained on a large
collections of text. Their capabilities include summarization,
question answering, classification, translation, and text generation.
These abilities may decrease the effort for research support.
Had to analyze and summarize massive data.

Language generation and the ability to retrieve information from evidence but not language generation.
Roles. An LLM is capable of generating comprehensible text even if not during training.
By the most recent scientific publications. This is a key point in the
knowledge-cutoff problem.

\subsection{Knowledge Cutoff}

Knowledge cut off is used here to mean a time cutoff point after which.
Relevant information is not represented in the available model
Parametric information. When a paper falls outside those boundaries and goes past this boundary and.
Because it can be inferred by the model if it is not provided by an outside source.

Knowledge cutoff is not the same as hallucination. Cutoff is
primarily an information-availability problem, whereas hallucination
Be intertwined, with misinformation potentially leadng to hallucinations, as a hallucinating model can hallucinate a misinformation response.
Interact: a fresh evidence despari talcan fill an information gap!
With older knowledge or unsupported statements.

\subsection{LLM-Assisted Literature Reviews}

An LLM-enabled literature review might do a search,
paper summarization, comparison of findings, thematic classification,
Evidence synthesis, and drafting. While these activities have varied.
Requirements. Summarizing a provided paper for the model, however, is the same
From requesting the newest research in an area.

For the experiments, the test and train data are labeled in the first and second sentence.
Will differ based on the task and dataset used. Khraisha et al. compared the performance of GPT-4 on a.
In conjunction with screening and data extraction we found that performance was dependent on
task conditions, dataset balance, and prompt quality
\cite{khraisha2024}. This suggests the need for approaching a LLM-assisted review as a
Multi-stage workflow rather than just one generation step.

\subsection{Fast Moving Areas}

The time issue is particularly relevant in AI,
Cybersecurity, medicine and biotechnology. In this spheres, new
models, methods, datasets, vulnerabilities, clinical findings, and
Benchmarks can mutate the research mosaic in a jiffy. An examination that was
Sufficient at a time may become short as evidence emerges.

\section {Thematic Literature Review}

\subsection{Recent Work Lacking}

At the most basic level a knowledge cutoff's most immediate consequences is the omission of studies.
Published beyond the available knowledge frontier. A review created
Without current retrieval could thus reflect a prior state of a
field.

Absent recent work can influence descriptive and analytical.
Sections. A subsequent study could establish a new benchmark, question an
If it's missing, the review can't provide earlier result, orit can't extend.a technique. If it's missing, the review cannot.
Incorporate that contribution.

\subsection{Outdated Information}

A model will sometimes correctly describe a technique that is relevant only for a number of years.
Roughly a decade ago without highlighting the newer options. This results in a.
Temporal mismatch with the review created and the present
research landscape.

This problem becomes particularly symptomatic when the user requests the „latest".
State of the art. Such a request requires explicit access to current
"Not based on evidence, not on insider knowledge, but evidence."

\subsection{Incorrect Research Gaps}

Identified gaps in the evidence include: The evidence is danged.
Lating the gaps in knowledge. In the absence of contemporary literature, an LLM
May refer to an issue as still unresolved if the subsequent research has notyet 
addressed it.

A statement like "a limited number of research has examined X" must, therefore, refer to.
A search that is both broad and up to date enough. Gap identification is not just
Not a language-generation task, but an evidence-coverage task.

\subsection{Temporal Bias}

Older publications may be simpler for a model to reproduce than very.
Recent literature. So, an LLM-generated review can provide.
Giving disproportionate focus to historical work Temporal bias can influence in which
research trends and apparent importance are interpreted.

A trustworthy review should then consider the distribution of.
Using evidence from across time and specifically looking for recent publications,
Rather than using refs created often, as proof of
current importance.

\subsection{Citation and Reference Errors}

The freshness of the knowledge is different from the reliability of the citation A citation
Might seem to be authentic as it has author, year, journal and DOI,
However, formatting does not prove that the publication is available or that
it supports the associated claim.

Walters and Wilder provided false and false bibliographies.
Sample of ChatGPT literature reviews with citations in the hundreds
\cite{walters2023}. Chelli et al. similarly evaluated references
Forschungsbericht arbeiten! For systematic reviews and had low precision and.
Substantial hallucination rates in the tested models
\cite{chelli2024}. These findings demonstrate why generated references
must be independently checked.

\section{Citation Reliability and Hallucination}

\subsection{Fabricated References}

An LLM might create a credit that looks like a authentic academic
Publication even if the publication does not exist. References such as this one:.
Can give the false appearance that unsupported statements have been made in support of evidence.

\subsection{Bibliographic Errors}

Even if a reference is an accurate publication, metadata may
Be wrong. Mistakes might appear in author names, titles, years, journal
Details, page information, or DOI identifications. For verification should be
check the complete bibliographic record.

\subsection{Citation-to-Claim Alignment}

Having a citation is not enough, the publication being cited must also
That is used to encourage the argument. A real paper can be used out of.
Grounds for a more confident conclusion than the paper justifies.

A practical verification process therefore checks:

\begin{enumerate}

\item whether the publication exists;

\item if the bibliographical metadata are valid; or.
\item determine if the publication does in fact support the specific claim.

\end{enumerate}

\subsection{Writing a Scientific Text}

Gao et al. demonstrated that ChatGPT has the ability to produce realistic scientific content
Abstracts that blinded human readers from generated and
Authentic abstracts somewhat hard to differentiate at time
\cite{gao2023abstracts}. It demonstrates the importance of the writing quality which should not
be treated as evidence of factual reliability.

Alkaissi and McFarlane also wrote of concerns with artificial
hallucinations in scientific writing \cite{alkaissi2023}. Together,
These studies require explicit validation of both the claims generated.
And references.

\section{Existing Approaches and Research Developments}

\subsection{Web and Scholarly Retrieval}

Search on the web can give an LLM access to information that is more recent than the model itself.
Knowledge. Scholarly databases add research-oriented filtering and
Bibliographic metadata. But search results do not include duplicates,
irrelevant material, non-peer-reviewed sources, or incomplete coverage.

Therefore, it is prudent to screen for post retrieval.
Relevance, publication quality, date, and methodological fit.

Search tools of research discovery can also assist researchers in finding related and.
Recently published studies through citation and similarity
relationships. Figure~\ref{fig:researchrabbit} shows an example of the
Research Rabbit interface for discovering related scholarly articles
And citation relationships. Tools like these can aid in the detection
You are receiving this message due to your readiness for the knowledge boundary of current research.
Of an LLM.

\begin{figure}[H]

\centering

\includegraphics[width=\columnwidth]{figures/researchrabbit.png}

\caption{Research Rabbit interface showing related scholarly articles
and citation relationships for literature discovery.}

\label{fig:researchrabbit}

\end{figure}

Retrieval-Augmented Generation (RAG) combines information retrieval
For Language Generation. Lewis et al. proposed a RAG framework that
combines parametric and non-parametric memory for knowledge-intensive
Tasks \cite{lewis2020rag}. Later a survey describes RAG as a way to
address problems including outdated knowledge by incorporating
information from external databases \cite{gao2023rag}.

RAG directly concerns knowledge cut-off as: current docs.
Can be retrieved at query time. While retrieval retrieves, can it guarantee
What a good review if the search documents are not related, are old,
Or duplicated, or poor quality generation can all still be inaccurate.

\subsection{ Reasoning and External Interaction}

ReAct marries reasoning trails and actions that permit a language
Model to connect to outside resources \cite{yao2023react}. For
Literature-review-help, this kind of interaction is important
Because searching, screening, and verification may need to be done...
Actions rather than a single response.

\subsection{Self-Reflective Retrieval}

As for the Self-RAG, it goes beyond retrieval-augmented generation by letting a model
Search for the information when needed and think about the passages you have obtained and.
Those of its designers \cite{asai2024selfrag}. It is a generalization to.
Synthesising information in different areas is important because accuracy of citation and verbatimness depend
Not only on document retrieval but also on assessment of their
usefulness.

\subsection{Corrective Retrieval}

Corrective RAG adds a retrieval evaluator that evaluates the retrieved
Documents and can also initiate supplementary retrieval operations.
Initial results are inadequate \cite{yan2024crag}. This direction
Shows that an important research trend emerging in high- quality retrieval is:
To be assessped to be correct.

\subsection{LLM-Assisted Literature Reviews}

An LLM-supported review might consist of: Searching assistance,
paper summarization, comparison of findings, thematic classification,
Evidence synthesis, and writing. All these activities have different
Requirements. Generating a summary of a paper been fed to model is a different
To turn to the newest research in a field.

Additional evidence from systematic-review tasks also indicates that LLM performance
Name varies by task and dataset. Khraisha et al. processed GPT-4 for the above experiment.
Choosing the relevant search criteria for screening and data extraction and discovered that the performance was based on:
task conditions, dataset balance, and prompt quality
\cite{khraisha2024}. This supports the argument to regard the LLM-aided review as a
Rather than a single generation step.

\subsection{Rapidly Evolving Fields}

This temporal problem is of particular importance in the field of AI.
Cybersecurity, medicine, and biotechnology. In such fields, new
models, methods, datasets, vulnerabilities, clinical findings, and
Why benchmarking can be the ultimate research killer A mild-mannered. A review that was.
Sufficient at one stage but may become full when further information becomes available.

\section{Thematic Literature Review}

\subsection{Missing Recent Research}

One of the most direct impact from a knowledge cutoff is missing of study
Published after the data knowledge boundary. A review conducted
Withholding current access to retrieval hmay therefore constitute an earlier state of a
field.

Missing current work can And descriptive and analytical
Sections. This more recent work can define a new standard, question a
Based on some previous finding, or are requesting to publicize this assessment. If this is not available, the standard is unable to.
Incorporate that contribution.

\subsection{Outdated Information}

A model can accurately represent a method that was significant for many.
Years ago while omitting newer options. This generates a
Temporal mismatch between the created review and the to run in the model:.
Research landscape.

This is especially true when the user has asked for the `latest''
State of the art. Such a request requires explicit access to current
Evidence and not static memos stored in mind.

\subsection{Wrong Research Gaps}

Research gaps emerge from the synthesis of evidence and
De in what is not clear. In the absence of current papers, an LLM
May identify an issue as open even if subsequent research has previously
addressed it.

A statement such as `few studies have explored X''must therefore
A wide and recent enough search. Not only is gap finding.
A language-generation task; it is an evidence-coverage task.

\subsection{Temporal Bias}

Older publications could be more subject to model replication than very.
Recent writings. As a result, an LLM summary can provide an entire overview of.
Attention is paid to earlier work. Temporal bias can influence.
Research trends and apparent importance are interpreted.

A robust review ought to, therefore, take a look at the distribution of
Look for evidence over time and specifically exclude the most recent articles,
Instead of considering commonly produced references as evidence of
current importance.

\subsection{Citation and Reference Problems}

-- Citation accuracy is a matter different from how current the information in the citation is. A citation
May appear believable as there is an author, year, journal, and DOI.
But that formatting does not prove the publication exists or that
it supports the associated claim.

Walters and Wilder cited in the work of book and illiterate.
Citations in a large dataset of ChatGPT authored literature reviews
\cite{walters2023}. Chelli et al. similarly evaluated references
Produced for systematic reviews and performance center!
Substantial hallucination rates in the tested models
\cite{chelli2024}. These findings demonstrate why generated references
Must be verified separately. Eval makes it easier to get up-to-date evidence, but it creates a new.
Failure reason: retrieval success. Confirmation of the reference and
Source accuracy, but one cannot recover citations that have not been obtained.
Problems with interpretation, but this can add to
Time and effort. Here they also highlight why hybrid workflows have been an
important research direction.

\section{Research Trends and Developments}

The literature tracks a development from static generation to the following:
Systems linking language models to external knowledge. Early.
Retrieval-augmented generation combined parametric and non-parametric
memory \cite{lewis2020rag}. Subsequent work has emphasized adaptive
Retrieval, self-correction and correction of inappropriate retrieval results
\cite{asai2024selfrag, yan2024crag}.

The other thing, and this is probably a function of where I am, is more emphasis on citation and factuality.
Analysis of the generated references shows that a fluent output can.
Contain fabricated or incorrect citations \cite{walters2023,chelli2024}.
This has converted the research problem from "can an LLM generate a
Can an LLM generate a review with the evidence can be A review?''

traced and checked?''

The third trend is human--AI task decomposition. Systematic evidence from
Review experiments indicate that LLMs can help with certain
Screening or extraction tasks in appropriate conditions, and
performance varies across datasets and tasks
\cite{khraisha2024}. This supports processes where automations can be.
Combined with human control rather than leapfrogging every step of the way.
Be automated equally well.

\section{Challenges and Limitations}

\subsection{Retrieval Coverage}

There is no retrieval system that is infallible.
Each database has a unique collection of journals and document types, and
Keyword-based searching will overlook papers using other words.

\subsection{Source Quality}

Even up-to-date information may be inacurate if the source you retrieved is
Poor quality Search freshness would therefore not be equated with
scientific validity.

\subsection{Conflicting Findings}

"Recent work may contradict earlier work. Each system must recognise"
Or if a paper is novel, corroborates, extends or refutes previous work
A different approach to prioritising rather than just giving all papers equal value.

\subsection{Citation Verification Cost}

It can become costly to verify every claim and cite that you generate. A
A practical system thus requires processes for ranking claims that
Are particularly significant, recent, controversial, or poorly documented.

\subsection{Human Interpretation}

Research gaps and scientific significance require interpretation.
Automated machines are capable of sorting evidence; but whether a
Represents a substantial gap in knowledge.
Domain-dependent.

\section{Research Gaps}

However, there is much work still to be done in retrieval and verification.

\subsection{Standardized Temporal Evaluation}

A requirement is for evaluation protocols that directly assess the

Well a generated review of all literature over time. So:

Use benchmarks independently evaluate the now coverage and the coverage
of recent publications.

\subsection{Temporal Retrieval Recall}

Additional tools should evaluate the retrieval recall for the nearly-studies in future systems.
Originality of source, quality of source, and source types, and temporal diversity. A review
Alone not be taken as current just because it has a few
recent citations.

\subsection{Citation Correctness}

Go beyond verification of whether a reference.
. Systems should also verify bibliographical accuracy and if
the cited publication supports the generated claim.

\subsection{Research-Gap Identification}

Automated identification of gaps is still challenging as it involves
Cross-paper comparison the focus on innovation and currency. Systems
Be considered when assessing if the suggested gaps are present.
Retrieving the current literature.

\subsection{Human--AI Collaboration}

Future work should explore how to apportion between researchers and LLMs.
Tasks effectively. Systems could automatically process routine
Rates while redirecting doubtful, high-stakes assertions to a human
reviewers.

\subsection{Evaluation in the Fast Moving Fields}

Lack of movement in many benchmarks domains that move quickly

Demand periodic reviews and that reflect the following:

new terminology, methods, benchmarks, and changing evidence.

\section{Suggestions for Future Research}

Future literature-review systems should provide explicit source
Adding linking so that significant assertions can be linked to evidence.
Documents. This would enable researchers to review the evidence instead of).
Trusting generated prose.

An alternative avenue is automated temporal analysis. Systems might
Study in order of publication and note if new studies are
confirms, extends, or contradicts earlier findings.

Citation verification should also become more claim-aware. Instead of
By validating the presence of a DOI using only that compares generated systems should.
Compares claims with the published paper and identifies potential disagreements.

Uncertainty indicators could further distinguish well-supported claims
From statements based on sparse or contradictory evidence. This
Would assist researchers target manual validation where it could be most
needed.

Finally, human-in-the-loop systems should be built on real
Research processes. Not everything is meant to be automated,
But well organized division of Retrieval and Generation,.
Verification, and expert judgment.

\section{Recommended Workflow}

From the reviewed literature, the following should be delineated for a reliable workflow:
retrieval, screening, evidence extraction, generation, and final
verification. Figure~\ref{fig:workflow} presents the proposed process.

\begin{figure}[t]

\centering

\begin{tikzpicture}[>=Latex,

step/.style={

draw,
 rounded corners,
 align=center,
 text width=32mm,

minimum height=7mm

}]

\node[step] (s1)

{1. Search recent academic references};

\node[step, below=3.5mm of s1] (s2)

{2. Retrieve and screen relevant papers};

\node[step, below=3.5mm of s2] (s3)

{3. Verify metadata and source quality};

\node[step, below=3.5mm of s3] (s4)

{4. Extract evidence and summarize};

\node[step, below=3.5mm of s4] (s5)

{5. Synthesize findings with the LLM};

\node[step, below=3.5mm of s5] (s6)

{6. Check claims and citations};

\node[step, below=3.5mm of s6] (s7)

{7. Human review and final interpretation};

\foreach \a/\b in {s1/s2,s2/s3,s3/s4,s4/s5,s5/s6,s6/s7}

\draw[->] (\a) -- (\b);

\end{tikzpicture}

\caption{Recommended workflow for reliable LLM-assisted literature reviews.}

\label{fig:workflow}

\end{figure}

As shown in Figure~\ref{fig:workflow}, retrieval occurs before
Synthesis to increase the amount of recent evidence that can enter into the workflow. Screening and.
Results with incorrect or irrelevant information.
The sources. Organization and synthesis with the LLM.
Claim and citation reviews are carried out in advance of the final human review.

This process has intentionally not been fully automated. Human in the loop.
Responsible for evaluating conflicting evidence, scientific
significance, and research-gap claims.

\section{Conclusion}

Knowledge cutoff is another critical limitation when LLMs are employed to
Generate literature reviews in rapidly evolving research fields. A
Model might not have seen very new papers, and may memorize them,
We do not identify research gaps accurately, or produce nonsensical citations.

Reviewed in this paper demonstrates that various mitigation
Strategies tackle various facets of the problem. Web and academic.
Retrieval get you the most up-to-date evidence; RAG combines retrieved
information with generation; self-reflective and corrective approaches
evaluate retrieval and generation; citation verification checks
– Bibliographical and claim-level reliability is better; and human review with backing.
Interpretation.

The key takeaway is that LLMs must be regarded as research.
Assistants less than the authorities on the current
Scientific papers. An up-to-date review has to make use of updated.
Retrieval, verification, structured synthesis, and human judgment.
In the future, research in a standardized time frame,
retrieval coverage, citation correctness, research-gap identification,
and effective human--AI collaboration.

\bibliographystyle{ACM-Reference-Format}

\bibliography{references}

\end{document}